\section{Limitations}
\label{sec:limitations}

\begin{itemize}
  \item \textbf{Scale and evaluation.} Models of 10M--160M parameters with a 16k vocabulary, far below the sizes at which
  looped language models are deployed, evaluated by validation loss only, without downstream benchmarks. The trends
  with size may bend; there are five sizes, one seed at 80M and one looped cell at 160M.
  \item \textbf{Data.} One source, FineWeb-Edu, with no document seen twice. Data-limited or repeated-data regimes, where loops may
  matter more \citep{chen2026growth,gulbahar2026squeezing}, are not covered. Seeds change only the initialisation,
  except in the data-order experiment.
  \item \textbf{Equal compute.} 62 of the 320 equal-compute comparisons rest on an extrapolated dense loss, bracketed by
  two methods (Section~\ref{sec:equalcompute}), and the measured comparisons are not independent of each other.
  \item \textbf{Learning rates.} The dense rates come from a sweep at one budget; the check of looped cells at 20M
  tried rates only up to the dense rate, which nine of its ten cells preferred; looped rates were re-tuned for one cell at
  10M--40M with one seed (two at 20M), and the other cells and 80M--160M use the dense rates, which understated the
  looped model by up to 0.024 nats at 10M, where at $80N$ and $160N$ its best rate was the lowest tried, so its optimum
  is not bracketed. On
  the synthetic task, the loops = steps and chain-of-thought models tied at the top of the learning-rate criterion and
  the lowest rate was taken, the 4-loop, random-loop and S9 models were best at the lowest rate tried, and extrapolation is
  sensitive to the rate.
  \item \textbf{Single seeds} at 80M, for the random-loop language models, and for the KV-sharing test.
  \item \textbf{One synthetic task family.} The serial-task evidence comes from chains of table lookups and 6.6M-parameter
  looped models; extrapolation length varies widely across seeds. A pointer-chasing control was not learned by the four
  conditions tried on it and was dropped.
  \item \textbf{Depth on text.} All dense references have eight layers and differ only in width, so the multiplier also
  contains whatever a deeper, narrower dense model of the same size would gain; we trained no such model, and no dense
  control of equal executed depth either.
  \item \textbf{The comparison with a chain of thought} covers one task, models of 6.5M--29M parameters and one training
  recipe. Its prompt grows with the number of steps, so the chain of thought's failure beyond the training range is a
  failure to generalise to longer inputs, which the looped models handle on the same prompts; whether a written chain
  extends on a task whose input length does not grow was not tested. The chain-of-thought models were supervised on every intermediate
  value, the looped models only on the answer; they got no help with length generalisation, and continuous thoughts \citep{hao2025training} were not tried.
  \item \textbf{Deployment} was measured on one GPU type, for models of at most 160M parameters, with one KV cache per
  executed layer; KV sharing was tested only as a change after training, at 20M.
  \item \textbf{Exploratory analyses.} The multiplier without a fitted law, the equal-compute bounds and the probes of the
  language models were chosen after seeing data and are labelled exploratory. The stopping comparison (S8) was planned
  after the frozen rule's results were known, though frozen before it ran, and its selection criterion picked the
  confidence rule at 0.99 that the text highlights only for $r = 8$ and the random loop count, and only once votes are
  set aside. Two readings of the frozen synthetic predictions, P3c's start at 10 steps and P4's diagonal share, were
  fixed after the results (Section~\ref{sec:predictions}). The pre-registered verdicts stand on their own.
\end{itemize}
